\subsection{An independent finite separation certificate}

\begin{claim}{A.50}{Proved}{A common cyclic-nineteen witness for the two characteristic-two points}
Let
\[
 C/\mathbf F _2:\quad y^2+x(x+1)y=x^5+1,
 \qquad P=(0,1),\quad Q=(1,0),\quad O=\infty.
\]
Set $L=\mathbf F _8=\mathbf F _2[t]/(t^3+t+1)$ and $J=J_{C_L}$.
Then $\#J(L)=76$ and $H=J(L)[2^\infty]$ has order four.
The Lang map $F_8-1:J\longrightarrow J$ factors as
\[
 J\longrightarrow J/H\xrightarrow{\lambda}J,
 \qquad \deg\lambda=19.
\]
Pulling $\lambda$ back along $j_P:R\mapsto[R-P]$ gives a connected,
constant cyclic, finite \emph{\'etale} cover $D\to C_L$ of degree
nineteen, with the chosen point $(P,0)$ above $P$. Its fibre above $Q$
also splits completely. It has genus twenty and
\[
 \frac{\#J_D(L)}{\#J(L)}=12284193340665686881,
\]
an odd integer. In particular this cover is a class-number witness at
both $P$ and $Q$, independent of the proof of the point-section theorem.
The same construction based at any of the nine points of $C(L)$ is a
class-number witness. Its relative class number is the displayed integer
at $P,Q,O$ and is $81428104953590501$ at each of the other six points.
\end{claim}

\begin{proof}
We give the finite data, their interpretation, and the collision criterion.
First we check the geometric model. The partial derivative with respect
to $y$ is $x(x+1)$. At its zeros the derivative with respect to $x$
is $y+x^4$, which equals one at both points on the curve. Thus the
affine model is smooth. Dividing by $x^2(x+1)^2$ gives the
Artin--Schreier equation for $y/(x(x+1))$. Its poles at $1$ and
infinity are simple; at zero the pole of order two becomes simple
after subtracting $(1/x)^2+1/x$. Thus there are precisely three
simple poles. Their different exponents are two, and
Riemann--Hurwitz gives genus two. There is a unique ramified point
at infinity. The degree-two hyperelliptic pencil has canonical class
$2O$; equivalently the regular differential $dx/(x(x+1))$ has divisor
$2O$.
Write an element of $L$ as the integer whose binary digits are its
coefficients in $1,t,t^2$. In that convention the points, in the order
used below, are
\[
 (0,1),(1,0),(3,2),(3,4),(5,4),(5,6),(7,2),(7,6),O.
\]
Direct enumeration over $\mathbf F _2,\mathbf F _4,\mathbf F _8,
\mathbf F _{64}$ gives respectively $3,3,9,87$ points. The genus-two
identity
\[
 \#J(\mathbf F _q)
       =\tfrac12\bigl(\#C(\mathbf F _q)^2+
                           \#C(\mathbf F _{q^2})\bigr)-q
\]
follows from the first two Newton identities and the reciprocal Weil
polynomial. At $q=8$ it gives $76$. An independent enumeration of all
reduced Mumford representatives gives exactly seventy-six as well.
Here a representative is $(u,v)$ with $u$ monic of degree at most two,
$\deg v<\deg u$, and $u\mid v^2+(x^2+x)v+x^5+1$.
The program gives every representative and the whole addition table;
it checks all $76^3$ associativity triples and all inverses.

We specify the arithmetic sufficiently to reproduce that table.
For $D_i=(u_i,v_i)$ put
\[
 d=\gcd(u_1,u_2,v_1+v_2+h),\qquad
 s_1u_1+s_2u_2+s_3(v_1+v_2+h)=d,
 \quad h=x^2+x,
\]
and form
\[
 u=\frac{u_1u_2}{d^2},\qquad
 v=\frac{s_1u_1v_2+s_2u_2v_1+s_3(v_1v_2+f)}d\pmod u,
 \quad f=x^5+1.
\]
While $\deg u>2$, replace $u$ by the monic associate of
$(f+hv+v^2)/u$ and $v$ by $h+v$ modulo this new polynomial.
These formulas are the ideal product and reduction for the quadratic
algebra $L[x,y]/(y^2+hy-f)$: cancellation of the common ideal is
encoded by the displayed Bezout relation, and replacing the ideal by
its complementary factor in the norm gives the reduction formula.
All divisions have zero remainder in the retained execution.

Multiplication by four maps the seventy-six representatives onto a
cyclic group of order nineteen with kernel $H$. Choose the generator
specified in the JSON certificate; taking discrete logarithms after
multiplication by four gives a homomorphism
$\chi:J(L)\twoheadrightarrow\mathbf Z/19$. Its table is checked on
all $76^2$ pairs. On the displayed points its values on $[R-P]$ are
\[
       0, 0, 1, 18, 7, 12, 8, 11, 0.                 \tag{FC.1}
\]
The same list holds with $Q$ or $O$ as base point. Direct Mumford
arithmetic checks that $[P-Q]$ is nonzero of exact order two.
Geometrically, $\operatorname{div}(x/(x+1))=2P-2Q$, since the two
points are ramification points of the degree-two $x$-map; a principal
divisor $P-Q$ would give a degree-one map to the projective line,
contradicting the genus. This also proves its exact order without
relying on the table.

The Lang map is separable, with constant kernel $J(L)$. Quotienting its
source by $H$ leaves a separable isogeny with constant kernel
$J(L)/H$ of order nineteen. The Abel map induces the isomorphism on
prime-to-characteristic geometric abelian fundamental groups established
in the Kummer compatibility argument. Its pullback is therefore a
geometrically connected cyclic cover. The point $P$ splits by
construction. Equation (FC.1) shows that $Q$ splits too. Equivalently,
the Frobenius translation on its fibre is $\chi([Q-P])=0$.
Riemann--Hurwitz gives $g(D)=1+19(g(C)-1)=20$.

For completeness we compute the relative class number, not only its
parity. For a nontrivial character $\zeta^\chi$ of the cyclic deck
group, expand the unramified Euler product by effective divisors.
The resulting polynomial in genus two is
\[
 L_P(\chi,T)=1+
   \left(\sum_{R\in C(L)}\zeta^{\chi([R-P])}\right)T
       +8\zeta^{\chi([K_C-2P])}T^2.
\]
Indeed Riemann--Roch makes every coefficient above degree two vanish
and gives the displayed coefficient of degree two; the sum of a
nontrivial character over $J(L)$ is zero. The canonical class is
$2O=2P$, so the last root of unity is one. Therefore
\[
 L_P(\chi,1)=12+\zeta+\zeta^{18}+\zeta^7+
                      \zeta^{12}+\zeta^8+\zeta^{11}.
\]
Factoring the regular-representation Euler factor at each closed point
gives $Z(D,T)/Z(C,T)=\prod_{i=1}^{18}L_P(\chi^i,T)$. Evaluation at
one gives the asserted class-number ratio. The independent integer
resultant is
\[
 \operatorname{Res}_Z\left(\sum_{j=0}^{18}Z^j,
      12+Z+Z^{18}+Z^7+Z^{12}+Z^8+Z^{11}\right)
           =12284193340665686881.
\]
There is also a short parity check: $2$ has order eighteen modulo
nineteen, so the first polynomial is irreducible over $\mathbf F _2$;
the second polynomial reduces to six nonzero monomials and is not
its multiple. For the other base points, subtract the corresponding
entry of (FC.1) from every exponent and use
$\chi([K_C-2R])=-2\chi([R-O])$. The same finite resultant calculation
gives the second integer in the statement. Each is odd.

We finish by proving why an odd relative class number excludes a
collision, without invoking A1. If the two sections of $C_L$ were
geometrically conjugate, choose an invariant geometric path between
their base points. Its lift to the pointed cyclic cover produces an
$L$-rational $Q'$ above $Q$ and geometrically conjugate sections at
$P',Q'$. The Kummer compatibility applied on $D$ implies
$T=[Q'-P']\in J_D(L)[2^\infty]$. The class is nonzero because the
Abel map of a positive-genus smooth proper curve is injective on
geometric points.

Let $B=(\ker(J_D\to J))^0$. Pullback followed by norm is
multiplication by nineteen, and $J\times B\to J_D$ is an isogeny.
Thus $\#B(L)$ is exactly the computed odd ratio. On the rational
two-primary group the averaging idempotent
$19^{-1}\sum_{g\in\mathbf Z/19}g$ has complementary image in
$B(L)[2^\infty]=0$. Hence every deck transformation fixes $T$.
For a generator $g$ this gives the equivalence of the distinct disjoint
effective divisors $gQ'+P'$ and $Q'+gP'$. They define a degree-two
pencil. Such a pencil on a curve of positive genus in characteristic
two is separable: a purely inseparable degree-two map to the projective
line would identify the square subfield with a rational function
field and force genus zero. Thus $D$ is hyperelliptic. Its unique
hyperelliptic involution is central in its automorphism group. The
order-nineteen deck transformation descends to an order-nineteen
automorphism of the projective line. A fixed point of this latter
automorphism has at most two geometric points above it, and an
odd-order permutation fixes both. This contradicts the free action
of the nonidentity deck transformation on an \emph{\'etale} cover.
All implications used in this finite test have now been proved.
\end{proof}

\paragraph{Check record.}
Checked only in this writing: GPT-6 Astra Pro, finite-checks reader,
18 November 2026. We wrote independent finite-field and divisor-class
arithmetic, enumerated the full group, checked the full operation and
character tables, and computed integer resultants. A second program
recomputed every addition using plain integer polynomial arithmetic
and all nine norms by integer Bareiss determinants. The effort setting
is not separately exposed by the runtime. Inputs, outputs, group
tables and cover definitions are in
\path{code/s2_f8_witness.py} and \path{code/s2_f8_witness.json}.

\subsection{A finite nonnilpotent test of the relative configuration sequence}

\begin{claim}{A.51}{Proved}{An explicit nonnilpotent cusp quotient and a dominating configuration quotient}
Let $\Gamma=P_2(S_2)$ be the pure braid group on two points of a
closed oriented genus-two surface, and let
$\Lambda=\langle a,b,c,d\rangle$ be its first fibre. For $p=2$,
$r=5$, and $A=\mathbf Z/3$, the fibre has a quotient
\[
 Q=\mathbf F _5^2\rtimes\mathbf Z/3,
 \qquad \sigma|_{\mathbf F _5^2}=
 S=\begin{pmatrix}0&-1\\1&-1\end{pmatrix},
 \qquad |Q|=75,
\]
in which the normal cusp subgroup is $\mathbf F _5^2$ and is
noncentral. An explicitly specified finite quotient $T$ of $\Gamma$
has order $9\cdot5^{40}$ and a normal exponent-five subgroup of
order $5^{40}$ and nilpotency class three with quotient $A\times A$. Its
restriction to $\Lambda$ dominates $Q$. Consequently the two maps
factor through the respective cuspidally pro-five groups over their
pro-odd proper quotients. This is a nonnilpotent finite test of the
relative configuration sequence.
\end{claim}

\begin{proof}
We use Bellingeri's Theorem 1.2 \citep[Theorem~1.2]{Bellingeri2004} for a closed oriented surface of
genus $g\ge1$, specialized to $g=n=2$. This gives generators
$a,b,c,d,\sigma$ for the full braid group. Its four same-generator,
four different-handle, two mixed same-handle, and one surface
relations are transcribed in the input certificate. Put
$t=[a,b^{-1}][c,d^{-1}]=\sigma^2$ and
$A_i=\sigma x_i\sigma^{-1}$, where
$(x_1,x_2,x_3,x_4)=(a,b,c,d)$. Index-two Reidemeister--Schreier
rewriting with representatives $1,\sigma$ gives a presentation of
$\Gamma$ with the eight generators $x_i,A_i$. Rewriting each of
the eleven relations and its $\sigma$-conjugate gives twenty-two
relators. During the rewriting, a positive $\sigma$-edge from the
second coset contributes $t$, a negative $\sigma$-edge from the
first contributes $t^{-1}$, and the other two edges contribute
one. This completely specifies the presentation computation.

Let $K$ be the kernel of $\Lambda\to\mathbf Z/3$ given by
$a\mapsto1$ and $b,c,d\mapsto0$. The cover of the compact surface
is the degree-three cover $C'\to C$; its punctured fibre has three
punctures and free group $K$ of rank ten. Representatives
$1,a,a^2$ give a free basis by all positively oriented Cayley edges
except the first two $a$-edges. The program retains the ten resulting
free words. Define
\[
 P=K/\bigl(\gamma_3(K)K^5\bigr).
\]
This is the free class-two exponent-five group on ten generators,
of order $5^{55}$. A concrete model is
\[
 P=V\oplus\bigwedge^2V,
 \quad (v,z)(w,u)=(v+w,z+u+\tfrac12v\wedge w),
 \quad V=\mathbf F _5^{10}.
\]
Every automorphism used below is specified by two matrices $(M,B)$,
of sizes $10\times10$ and $45\times10$, respectively, acting by
\[
              (v,z)\longmapsto(Mv,(\bigwedge^2M)z+Bv).
                                                               \tag{FC.2}
\]
The eight pairs of matrices in the certificate thus specify elements
of a concrete finite matrix group, not formal symbols.

For the first four generators, use the action on $K$ by conjugation
with $x_i$. For the second four, the point-push words are as follows:
write $u_i=t x_i^{-1}t x_i t^{-1}$ and $w_i=u_i^{-1}t$.
They send $x_i$ to $t x_i t^{-1}$. If $i$ is odd, the partner
$x_{i+1}$ goes to $t x_{i+1}t^{-1}u_i$; if $i$ is even,
$x_{i-1}$ goes to $x_{i-1}t^{-1}$. Earlier handles are fixed,
and later handles are conjugated by $w_i$. These formulas follow
by rewriting Bellingeri's same-handle and different-handle relations.
For an independently checkable safeguard, we do not rely on that
derivation alone: substituting the eight finite matrices into all
twenty-two relators gives the identity, exactly over $\mathbf F _5$.
Thus they define a homomorphism
$\rho:\Gamma\to\operatorname{Aut}(P)$ even independently of
any statement about automorphisms of a free group.

Append the two deck coordinates
$\delta:\Gamma\twoheadrightarrow(\mathbf Z/3)^2$, assigning
$(1,0)$ to $a$, $(0,1)$ to $A_1$, and zero to the other six
generators. Put $T=\operatorname{im}(\rho,\delta)$.
This defines a finite group explicitly as the subgroup generated
by eight specified matrices with their specified deck coordinates.
Let $H_0=\ker\delta$. Geometrically it is the fundamental group of
$C'\times C'$ minus the three deck graphs. Representatives
$a^iA_1^j$ for $0\le i,j<3$ give seventy-two Schreier generators
of $H_0$. Their linear parts are all $1+N_h$. If $W\subset V$
is the two-dimensional space spanned by the three cusp classes,
the certificate checks
\[
 N_h(V)\subset W,\qquad N_h(W)=0
 \quad\hbox{for all seventy-two generators}.
\]
It checks all $72^2$ identities $N_hN_{h'}=0$ as well; the span of
the $N_h$ has dimension four. Therefore the linear image of $H_0$
is an elementary abelian five-group. The kernel of
$\operatorname{Aut}(P)\to\operatorname{GL}(V)$ in (FC.2) is
$\operatorname{Hom}(V,\bigwedge^2V)$, also a five-group.
It follows that $\rho(H_0)$ is a five-group. We now determine its
order, without enumerating all of its elements.

Here is the finite logarithm argument used for this purpose. If $I$
is an associative algebra over $\mathbf F _5$ with $I^5=0$, the
polynomials
\[
 \log(1+D)=D-\frac{D^2}{2}+\frac{D^3}{3}-\frac{D^4}{4},
 \qquad
 \exp(X)=1+X+\frac{X^2}{2}+\frac{X^3}{6}+\frac{X^4}{24}
\]
are inverse. Multiplication is transported to the BCH polynomial
through bracket degree four; its denominators divide $24$, so they
are invertible in the field. For matrices $g_i=\exp(X_i)$, their
generated group is exactly $\exp(L)$, where $L$ is the Lie algebra
generated by the $X_i$. To prove the reverse inclusion, which is the
nontrivial direction, first work in the free nilpotent Lie algebra
of class four on the symbols $X_i$. Give a bracket word its ordinary
bracket length. The logarithm of an iterated group commutator equals
its corresponding Lie bracket modulo larger bracket lengths.
Integer powers supply all five scalar multiples. Products add the
terms of the current smallest length modulo larger lengths.
Consequently every homogeneous linear combination of bracket length
$j$ can be realized modulo terms of length $j+1$. Starting at $j=1$
and correcting successively through $j=4$ realizes the exponential
of every Lie element exactly. Mapping this free algebra to $I$
proves the assertion for the matrices in question. In particular,
$|\langle g_i\rangle|=5^{\dim L}$. The same argument identifies
successive lower central factors on the group and Lie sides.

To check the hypothesis, put $E_h=\bigwedge^2(1+N_h)$. The
operator $E_h-1$ lowers the filtration
\[
 \bigwedge^2 V\supset W\wedge V\supset\bigwedge^2 W\supset0
\]
by one step. The operators $N_h$ lower $V\supset W\supset0$
by one step. The full augmentation matrices have blocks
$\left(\begin{smallmatrix}N_h&0\\ B_h&E_h-1\end{smallmatrix}\right)$.
A product of five such matrices is zero: a nonzero block path uses
at most one $N$ step, at most one $B$ step, and at most two $E-1$
steps. Thus the algebra $I$ they generate satisfies $I^5=0$.

The certificate computes the seventy-two logarithms, whose linear
span has dimension twenty, and closes their span under brackets.
The resulting Lie algebra has dimension forty. Its lower central
series has dimensions
\[
                   40,\quad24,\quad4,\quad0.
\]
The retained basis and all bracket reductions are exact over
$\mathbf F _5$. Hence $\rho(H_0)$ has order $5^{40}$, exponent
five, and nilpotency class exactly three. Since the appended deck
coordinates are surjective,
\[
             |T|=9\cdot5^{40}
                =81854523159563541412353515625.
\]
A second computation uses no logarithms. It constructs a common
invariant flag of dimensions $19,41,47,55$, makes all seventy-two
matrices unitriangular, and collects actual matrix commutators in
increasing triangular pivot order. It obtains forty generators of
order five. All pair commutators collect, as do all original
generators. Collection terminates because every commutator raises
the triangular pivot filtration. Distinct ordered products have
distinct first unequal pivot coordinates, so their $5^{40}$ normal
forms are distinct. This independently confirms the order.

This construction also realizes the finite-core step of the relative
completion proof. The subgroup
$J_0=\ker(H_0\to\operatorname{Aut}(P))$ is invariant under
$\Gamma$, since it is the intersection of $H_0$ with a homomorphism
kernel. Thus its core over the nine cosets is already $J_0$.
The subgroup $K$ is preserved by every generator, as the displayed
words preserve its abelian degree-three quotient. Its quotient $P$
is characteristic. Hence the core just used has precisely the
cover, shear, and normality interpretation required by that proof.

It remains to check domination of the nonnilpotent fibre quotient.
Define $q:\Lambda\to Q$ by
\[
 a\mapsto(0,\sigma),\qquad b\mapsto((1,0),1),
 \qquad c,d\mapsto1.
\]
The matrix $S$ has order three and no invariant line over
$\mathbf F _5$. The images generate all seventy-five elements of
$Q$. The nonzero image of $t$ and its $S$-translate span the
normal subgroup $\mathbf F _5^2$, so this is exactly the normal
cusp subgroup. In particular $Q$ is nonnilpotent. The restriction
of $q$ to $K$ factors through a rank-two linear map
$\lambda:V\twoheadrightarrow\mathbf F _5^2$; its matrix is
provided in the certificate.

Suppose $x\in\Lambda$ is killed by $(\rho,\delta)$. Its first
deck coordinate shows that $x\in K$. Its action on $P$ is the
inner action of its image. The centre of $P$ is exactly
$\bigwedge^2V$: if $v\ne0$, choose $w$ independent of $v$, so
$v\wedge w\ne0$. Equivalently the certificate computes rank ten
for $v\mapsto(w\mapsto v\wedge w)$. Thus $x$ has zero image in
$V$ and is killed by $\lambda$, hence by $q$. This proves the
required kernel inclusion and therefore domination.

Finally the kernel of the proper-product quotient maps into
$\rho(H_0)$, a five-group. The maps therefore kill the kernels of
the maximal pro-five quotients of the corresponding cusp groups.
This proves the asserted factorization. No nilpotent quotient of
the entire fibre has been substituted for $Q$.
\end{proof}

\paragraph{Check record.}
Checked only in this writing: GPT-6 Astra Pro, finite-checks reader,
18 November 2026, source inspection, explicit presentation rewriting,
matrix computation, and the complete kernel argument above.
The effort setting is not separately exposed by the runtime.
A second program independently recomputed all eighty generator images
in a truncated tensor algebra and verified all twenty-two relations
as full $55\times55$ matrices. The source, all words, all matrices,
all shear tests and the factor map
are retained in \path{code/s3_nonnilpotent.py} and its JSON output.
The two independent order methods and their full bases are in
\path{code/s3_group_order.py} and \path{code/s3_order_collect.py},
with their JSON certificates.

\subsection{Independent reruns of the four retained finite claims}

\begin{claim}{A.52}{Proved}{Agreement of the independent finite certificates}
The finite assertions in A.7, A.12, A.27, and A.29 have been checked
by a newly written program from their mathematical inputs, before
opening their original program sources, and by rerunning those original
programs. The two computations agree. More precisely:
\begin{enumerate}
\item The two integral rank-eighty-two homology matrices of A.7 are
integrally conjugate by a supplied matrix of determinant $\pm1$.
Their restricted matrices are
$\left(\begin{smallmatrix}1&3\\0&1\end{smallmatrix}\right)$ and
$\left(\begin{smallmatrix}1&0\\-3&1\end{smallmatrix}\right)$,
and their product has the claimed nonzero minus-one eigenvector
modulo five.
\item The five transvections in A.12 agree entry by entry, generate
$51840$ matrices, and the specified word has exact order five.
\item The finite-field, point, orbit, and group-order calculations
in A.27 agree: $129$ points over $\mathbf F _{128}$, $126$ of exact
degree seven, and $|\operatorname{GL}_4(\mathbf F _3)|=24261120$.
\item The group in A.29 has order $576$, the stated generators,
commutator subgroup of order $64$, trivial complement core, and the
specified order-seven automorphism. All $576^2$ product
compatibilities hold.
\end{enumerate}
This statement concerns the finite certificates. The profinite and
geometric consequences still require the separate proofs in the
corresponding statements.
\end{claim}

\begin{proof}
For A.7 we computed the induced action on the graph-cycle space of
the Heisenberg cover using a breadth-first tree that uses positive
generators only. Its eighty-two non-tree edges give coordinates.
Substitution of the point-push words in each closed lifted word gives
the integral action. The original program uses a different, signed
tree. We therefore evaluated every original basis word in our basis,
obtaining a matrix $U$ with determinant $\pm1$, and verified
$M_{
m new}U=UM_{
m old}$ for both complete matrices. This
explains the only coordinate differences; the mathematical data agree.

For A.12 the independent computation starts from the geometric
transvection formula on $\mathbf F _3^4$, with directions
$e_1,e_2,e_3,e_4,e_1+e_3$ and coefficient two. Exhaustive multiplication
gives $51840$ distinct matrices and five successive powers verify
the exact order of the indicated word. The original program was
separately rerun to check its derivation from full Schreier homology.
We do not describe the independent calculation as a second derivation
of all of its free-group words.

For A.27 we implemented polynomial-basis binary-field arithmetic,
verified $a^{127}=1$ for every nonzero element of the proposed
field, enumerated its curve points, and computed the full seven-point
Frobenius orbit. For A.29 we independently enumerated all elements
of $\mathbf F _{64}\rtimes\mathbf Z/9$, all products required by
the automorphism test, the generator orbit and the conjugate
complement intersection. The complete outputs and comparisons are
retained. All assertions are finite identities or exhaustive finite
enumerations, so their retained programs constitute reproducible
certificates for the assertions just listed.
\end{proof}

\paragraph{Check record.}
Also checked in this writing, to the finite scope just stated:
GPT-6 Astra Pro, finite-checks reader, 18 November 2026. The blind
program hash, the precise scope of independence, the unmodified
rerun inputs, all outputs, and the comparison are recorded in
\path{S6_BLIND_PROTOCOL.json}, \path{code/independent_s6.py},
and \path{code/s6_comparison.json}.
